% Страницы 83--86. Текст и разметка восстановлены по предоставленному PDF.
Если работа внешних сил положительна, то кинетическая энергия увеличивается ($\Delta K>0$, тело разгоняется).

Если работа внешних сил отрицательна, то кинетическая энергия уменьшается ($\Delta K<0$, тело замедляет движение). Пример~--- торможение под действием силы трения, работа ко\-торой отрицательна.

Если же работа внешних сил равна нулю, то кинетическая энергия тела за это время не ме\-няется. Нетривиальный пример~--- равномерное движение по окружности, совершаемое грузом на нити в горизонтальной плоскости. Сила тяже\-сти, сила реакции опоры и сила натяжения нити всегда перпендикулярны скорости, и работа каждой из этих сил равна нулю в течение любого промежутка времени. Соответственно, кинетическая энергия груза (а значит, и его скорость) остаётся постоянной в процессе движе\-ния.

\vspace{10.331pt}\noindent
\textbf{Задача.} Автомобиль едет по горизонтальной дороге со скоростью $v$ и начинает резко тормо\-зить. Найти путь $s$, пройденный автомобилем до полной остановки, если коэффициент трения шин о дорогу равен $\mu$.

\vspace{5.165pt}\noindent
\textit{Решение.} Начальная кинетическая энергия автомобиля $K_1=\dfrac{mv^2}{2}$, конечная кинетическая энергия $K_2=0$. Изменение кинетической энер\-гии $\Delta K=K_2-K_1=-\dfrac{mv^2}{2}$.

На автомобиль действуют сила тяже\-сти $m\vec g$, реакция опоры $\vec N$ и сила трения $\vec f$. Сила тяже\-сти и реакция опоры, будучи перпендикулярны перемещению автомобиля, работы не соверша\-ют. Работа силы трения:
\[
A=-fs=-\mu Ns=-\mu mgs.
\]

Из теоремы о кинетической энер\-гии теперь получаем:
\[
\Delta K=A\Rightarrow-\frac{mv^2}{2}=-\mu mgs\Rightarrow s=\frac{v^2}{2\mu g}\,.
\]
\subsection*{1.16.5\quadПотенциальная энергия тела вблизи поверхности Земли}
\noindent Рассмотрим тело массы $m$, находящееся на некоторой высоте над поверхностью Земли. Высоту считаем много меньше земного радиуса. Изменением силы тяже\-сти в процессе перемещения тела пренебрегаем.

\begin{wrapfigure}{r}{189.306bp}
\vspace{-2.501pt}
\centering\hypertarget{fig150}{}\ReplicaGravityDiagram
\par\vspace{10.05pt}\hspace*{3.916bp}Рис.~1.50. $A=mg(h_1-h_2)$
\vspace{-6pt}
\end{wrapfigure}
Если тело находится на высоте $h$, то \textit{потенциаль\-ная энергия} тела по определению равна:
\[
W=mgh,
\]
\noindent где $g$~--- ускорение свободного падения вблизи поверхности Земли.

Высоту не обязательно отсчитывать от поверхности Земли. Как мы увидим ниже (формулы (\hyperlink{eq178}{1.78}), (\hyperlink{eq179}{1.79})), фи\-зическим смыслом обладает не сама по себе потенциаль\-ная энергия, но её \textit{изменение}. А изменение потенциаль\-ной энер\-гии не зависит от уровня отсчёта. Выбор нуле\-вого уровня потенциаль\-ной энер\-гии в конкретной задаче диктуется исключительно соображениями удобства.

Найдём работу, совершаемую силой тяже\-сти при перемещении тела. Предположим, что тело перемещается по прямой из точки $P$, находящейся на высоте $h_1$, в точку $Q$, находящуюся на высоте $h_2$ (рис.~\hyperlink{fig150}{1.50}).
\newpage
Угол между силой тяже\-сти $m\vec g$ и перемещением тела $\vec s$ обозначим $\alpha$. Для работы силы тяже\-сти получим:
\[
A=m\vec g\vec s=mgs\cos\alpha.
\]
\noindent Но, как видно из рис.~\hyperlink{fig150}{1.50}, $s\cos\alpha=h_1-h_2$. Поэтому
\[
A=mg(h_1-h_2)=mgh_1-mgh_2,
\]
\noindent или
\begin{equation}\hypertarget{eq178}{}A=W_1-W_2.\tag{1.78}\end{equation}
\noindent Учитывая, что $W_1-W_2=-(W_2-W_1)=-\Delta W$, имеем также:
\begin{equation}\hypertarget{eq179}{}A=-\Delta W.\tag{1.79}\end{equation}

Можно доказать, что формулы (\hyperlink{eq178}{1.78}) и (\hyperlink{eq179}{1.79}) справедливы \textit{для любой тра\-ек\-тории}, по ко\-торой тело перемещается из точки $P$ в точку $Q$, а не только для прямолинейного отрезка.

\vspace{5.165pt}
\textit{Работа силы тяже\-сти не зависит от формы тра\-ек\-тории, по ко\-торой перемещается те\-ло, и равна разности значений потенциаль\-ной энер\-гии в начальной и конечной точках тра\-ек\-тории. Иными словами, работа силы тяже\-сти всегда равна изменению потенциаль\-ной энер\-гии с противоположным знаком.}

\textit{В частности, работа силы тяже\-сти по любому замкнутому пути равна нулю.}

\vspace{5.165pt}
Сила называется \textit{консервативной}, если при перемещении тела работа этой силы не зависит от формы тра\-ек\-тории, а определяется только начальным и конечным положением тела. Сила тяже\-сти, таким образом, является консервативной. Работа консервативной силы по любому замкнутому пути равна нулю. Только в случае консервативной силы возможно ввести такую величину, как потенциаль\-ная энергия.

\subsection*{1.16.6\quadПотенциальная энергия деформированной пружины}
\noindent Рассмотрим пружину жёсткости $k$. Начальная деформация пружины равна $x_1$. Предположим, что пружина деформируется до некоторой конечной величины деформации $x_2$. Чему равна при этом работа силы упру\-гости пружины?

В данном случае силу на перемещение не умножишь, так как сила упру\-гости ме\-няется в процессе деформации пружины. Для нахождения работы переменной силы требуется интегри\-рование. Мы не будем приводить здесь вывод, а сразу выпишем конечный результат.

Оказывается, сила упру\-гости пружины также является консервативной. Её работа зависит лишь от величин $x_1$, $x_2$ и определяется формулой:
\[
A=\frac{kx_1^2}{2}-\frac{kx_2^2}{2}\,.
\]

Величина
\[
W=\frac{kx^2}{2}
\]
\noindent называется \textit{потенциаль\-ной энергией деформированной пружины} ($x$~--- величина деформации). Следовательно,
\[
A=W_1-W_2=-\Delta W,
\]
\noindent что полностью аналогично формулам (\hyperlink{eq178}{1.78}) и (\hyperlink{eq179}{1.79}).
\newpage
\subsection*{1.16.7\quadЗакон сохране\-ния механи\-ческой энер\-гии}
\noindent Консервативные силы называются так потому, что сохраняют механическую энергию замкну\-той си\-стемы тел.

\textit{Механическая энергия} $E$ тела равна сумме его кинетической и потенциаль\-ной энергий:
\[
E=K+W.
\]

Механическая энергия си\-стемы тел равна сумме их кинетических энергий и потенциаль\-ной энер\-гии их взаимодействия друг с другом.

Предположим, что тело совершает движение под действием силы тяже\-сти и/или силы упру\-гости пружины. Будем считать, что трения нет. Пусть в начальном положении кинетическая и потенциаль\-ная энер\-гии тела равны $K_1$ и $W_1$, в конечном положении~--- $K_2$ и $W_2$. Работу внешних сил при перемещении тела из начального положения в конечное обозначим $A$.

По теореме о кинетической энер\-гии:
\[
K_2-K_1=A.
\]

Но работа консервативных сил равна разности потенциальных энергий:
\[
A=W_1-W_2.
\]

Отсюда получаем:
\[
K_2-K_1=W_1-W_2,
\]
\noindent или
\[
K_1+W_1=K_2+W_2.
\]

Левая и правая части данного равенства представляют собой механическую энергию тела в начальном и конечном положении:
\[
E_1=E_2.
\]

Следовательно, \textit{при движении тела в поле силы тяже\-сти и/или на пружине механиче\-ская энергия тела остаётся неизменной при отсутствии трения.}

Справедливо и более общее утверждение.

\vspace{5.165pt}\noindent
\textbf{Закон сохране\-ния механи\-ческой энер\-гии.} Если в замкну\-той системе действуют только консервативные силы, то механиче\-ская энергия си\-стемы сохраняется.

\vspace{5.165pt}
При этих условиях могут происходить лишь превращения энер\-гии: из кинетической в по\-тенциальную и наоборот. Общий запас механи\-ческой энер\-гии си\-стемы остаётся постоянным.

\subsection*{1.16.8\quadЗакон изменения механи\-ческой энер\-гии}
\noindent Если между телами замкну\-той си\-стемы имеются силы сопротивления (сухое или вязкое тре\-ние), то механиче\-ская энергия си\-стемы будет уменьшаться. Так, автомобиль останавливается в результате торможения, колебания маятника постепенно затухают и т.\,д. Силы трения неконсер\-вативны: работа силы трения очевидным образом зависит от пути, по которому перемещается тело между данными точками. В частности, работа силы трения по замкнутому пути не равна нулю.

Снова рассмотрим движение тела в поле силы тяже\-сти и/или на пружине. Вдобавок на тело действует сила трения, которая за рассматриваемый промежуток времени совершает отрица\-тельную работу $A_{\text{тр}}$. Работу консервативных сил (тяже\-сти и упру\-гости) по-прежнему обозна\-чаем $A$.
\newpage
Изменение кинетической энер\-гии тела равно работе \textit{всех} внешних сил:
\[
K_2-K_1=A+A_{\text{тр}}.
\]

Но $A=W_1-W_2$, следовательно
\[
K_2-K_1=W_1-W_2+A_{\text{тр}}.
\]

Отсюда
\[
K_2+W_2-(K_1+W_1)=A_{\text{тр}},
\]
\noindent или
\[
E_2-E_1=A_{\text{тр}}.
\]

В левой части стоит величина $\Delta E=E_2-E_1$~--- изменение механи\-ческой энер\-гии тела:
\[
\Delta E=A_{\text{тр}}.
\]

Итак, \textit{при движении тела в поле силы тяже\-сти и/или на пружине изменение механи\-ческой энер\-гии тела равно работе силы трения.} Так как работа силы трения отрицательна, изменение механи\-ческой энер\-гии также отрицательно: механиче\-ская энергия убывает.

Справедливо и более общее утверждение.

\vspace{5.165pt}\noindent
\textbf{Закон изменения механи\-ческой энер\-гии.} Изменение механи\-ческой энер\-гии замкну\-той си\-стемы равно работе сил трения, действующих внутри си\-стемы.

\vspace{5.165pt}
Ясно, что закон сохране\-ния механи\-ческой энер\-гии является частным случаем данного утвер\-ждения.

Конечно, убыль механи\-ческой энер\-гии не противоречит общефизическому закону сохране\-ния энер\-гии. В данном случае механиче\-ская энергия превращается в энергию теплового движе\-ния частиц вещества и их по\-тенциальную энергию взаимодействия друг с другом, т.~е. переходит во \textit{внутреннюю энергию} тел си\-стемы.
